\chapter{Learning the Kernel, and What Convergence Costs}\label{ch:em-kernel}

Chapter~\ref{ch:em-method} established that the estimator is correct: the
E-step is exact, the surrogate ascends, and the recovery error of the
second-order coordinates falls approximately as
$\nseq^{-1/2}$ over the tested range. None of that says how long it takes,
and this chapter is about the difference. Correctness is a statement about the fixed point; what follows
is a statement about the \emph{map} that reaches it, and the two are
independent. An estimator can be exactly right in the limit and still be
routinely reported before it arrives.

The finding is that the map does not converge at one speed. The autoregressive
coefficient and the innovation \emph{shape} contract at rates that differ by a
factor of eight in the exponent, so a stopping rule calibrated on the first
certifies a kernel that has not converged in the second. Every non-Gaussian
claim in this thesis depends on the coordinate that settles last, which makes
this a prerequisite for Chapter~\ref{ch:em-results} and not a diagnostic beside
it. Section~\ref{sec:em-rate-theory} says analytically why two rates should be
expected and which coordinate should be slower; the rest of the chapter
measures both and reports where the prediction holds and where it does not.

\section{Why there are two rates}\label{sec:em-rate-theory}

\subsection{EM is a fixed-point map, and its Jacobian is missing information}

The M-step of Section~\ref{sec:em-algorithm} defines a map
$\theta^{(k+1)} = \mathcal M(\theta^{(k)})$, and a stationary point $\theta^\star$
of the marginal likelihood is a fixed point of it. Expanding about that point,
\begin{equation}
  \theta^{(k+1)} - \theta^\star
  \;=\; J(\theta^\star)\,\big(\theta^{(k)} - \theta^\star\big)
     \;+\; O\big(\|\theta^{(k)} - \theta^\star\|^2\big),
  \qquad
  J \;=\; \frac{\partial \mathcal M}{\partial \theta}\bigg|_{\theta^\star},
  \label{eq:em-fixedpoint}
\end{equation}
so EM converges \emph{linearly}, not quadratically, at a rate set by the
spectrum of $J$. The identification of that Jacobian is the central result of
\citet{dempster1977maximum}. All three informations below are evaluated at
$\theta^\star$ and conditioned on the same realised observation $x$:
$I_{\mathrm{com}} = -\E[\nabla^2 \log p_\theta(a,c,x) \mid x]$ is the
expected complete-data information given $x$,
$I_{\mathrm{obs}} = -\nabla^2 \mathcal L$ is the observed information, and
$I_{\mathrm{mis}}$ is their difference --- the information the latent
variables retain. With that convention, and $J$ written for a map acting on
column vectors,
\begin{equation}
  I_{\mathrm{com}} \;=\; I_{\mathrm{obs}} \;+\; I_{\mathrm{mis}},
  \qquad
  J(\theta^\star)
  \;=\; I_{\mathrm{com}}^{-1} I_{\mathrm{mis}}
  \;=\; \mathbb{I} - I_{\mathrm{com}}^{-1} I_{\mathrm{obs}} .
  \label{eq:em-missing-info}
\end{equation}
The rate of EM \emph{is} the fraction of the information that the latent
variables carry away. If nothing is missing, $J = 0$ and EM converges in one
step; as $I_{\mathrm{mis}} \to I_{\mathrm{com}}$ the largest eigenvalue of $J$
approaches one and progress per update approaches nothing.

One scope condition on \eqref{eq:em-missing-info} governs everything measured
in this chapter. The identity is the Jacobian of the \emph{exact} EM map,
whose M-step maximises $Q$ jointly. The mixture M-step of
Chapter~\ref{ch:em-method} does not: it cycles conditional maximisations and
stops after at most the sixteen inner sweeps of
Remark~\ref{rmk:em-inner}, which is ECM in the sense of
\citet{meng1993maximum} and generalised EM when the cycle is cut short. In the
quadratic idealisation with $r$ inner sweeps the implemented map has Jacobian
$G^{r} + (\mathbb{I} - G^{r})\,J$, with $G$ the block Gauss--Seidel matrix of
the conditional maximisations, which is \eqref{eq:em-missing-info} only in the
limit $r \to \infty$. So \eqref{eq:em-missing-info} is the benchmark this
chapter measures \emph{against}, not a formula for the map that ran, and a
contraction measured from a trace mixes missing-information slowness with
conditional-maximisation slowness. It is treated throughout as an ordering
heuristic for which coordinate should be slow, and never as a calibrated
prediction of how slow.

\subsection{From eigenvalues to iteration counts}

Equation \eqref{eq:em-fixedpoint} says a coordinate aligned with an eigenvector
of $J$ with eigenvalue $\lambda$ has its error contract as $\lambda^k$. Stopping
that coordinate when its relative distance to the limit first falls below a
tolerance $\tau$ therefore takes
\begin{equation}
  k(\tau, \lambda) \;\approx\; \frac{\log \tau}{\log \lambda},
  \qquad\text{so}\qquad
  \frac{k_2}{k_1}
  \;\approx\;
  \underbrace{\frac{\log \tau_2}{\log \tau_1}}_{\text{tolerance}}
  \;\cdot\;
  \underbrace{\frac{\log \lambda_1}{\log \lambda_2}}_{\text{rate}} .
  \label{eq:em-settle}
\end{equation}
Two coordinates of the same fit can thus require very different iteration
counts, and \eqref{eq:em-settle} separates the two reasons cleanly: one is a
choice of tolerance, the other is a property of the model. Two limits on the
formula are stated now and verified below. It is an \emph{asymptotic tail}
statement --- during the transient a trace can fall much faster than its
terminal slope, so \eqref{eq:em-settle} upper-bounds rather than predicts a
measured settling time. And it speaks about eigenmodes: a parameter
\emph{coordinate} generally mixes several modes of $J$, so a slope fitted to
one coordinate's trace is an empirical contraction factor for that
coordinate, identified with an eigenvalue only if a single mode dominates its
window. Section~\ref{sec:em-two-rates} measures the coordinate slopes
directly, under exactly those labels.

\subsection{The channel enters the identity twice}
\label{sec:em-two-effects}

Both effects of the observation channel are visible in
\eqref{eq:em-missing-info}, and they act on different terms.

\emph{On the rate.} Fitting on clean transition pairs leaves only the mixture
label $c$ latent. Passing the same chains through the channel makes the entire
state $a$ latent as well, so $I_{\mathrm{mis}}$ grows by everything the
posterior over $a$ does not resolve, $J$ moves towards the identity, and the
map slows. This is a statement about iterations and would persist at infinite
sample size. Two things it is not. It is a comparison of what is latent in the
two cases, not a claim that the clean case is fast in absolute terms: the
components of a mixture fitted to a Laplace innovation overlap heavily, label
information can then be a large share of $I_{\mathrm{com}}$ on its own, and no
measurement here separates label slowness from state slowness. And
Section~\ref{sec:em-information} compares the two arms at their endpoints and
by accuracy rather than by contraction rate, so it corroborates the direction
of this argument without isolating it as the mechanism.

\emph{On the accuracy.} The channel also shrinks $I_{\mathrm{obs}}$, and the
asymptotic covariance of the estimator is $I_{\mathrm{obs}}^{-1}/\nseq$. This
is a statement about data, not iterations, and it does not vanish by iterating
longer. Section~\ref{sec:em-information} measures it directly.

\subsection{Which coordinate should be slower, and why}
\label{sec:em-which-slower}

The model supplies the answer before any fit is run. The forward channel is
variance-preserving and the chain is normalised to unit marginal variance
($\ivar = 1 - \corr^2$, so $\kappa_2(A) = 1$): the second cumulant of the
noised variable is then preserved \emph{exactly} at every $t$, while the
fourth is damped,
\begin{equation}
  \kappa_2(X_t) \;=\; \kappa_2(A) \quad \text{for all } t,
  \qquad
  \kappa_4(X_t) \;=\; e^{-4t}\,\kappa_4(A).
  \label{eq:em-cumulant-damping}
\end{equation}
The autoregressive coefficient is a second-order quantity: it is identified by
the lag-one covariance, which the channel attenuates in a known, invertible way
and does not destroy. The innovation \emph{shape} is a fourth-order quantity,
and \eqref{eq:em-cumulant-damping} says the channel removes an
$e^{-4t}$ fraction of exactly the structure that identifies it. The information
the channel takes away is therefore concentrated on the shape coordinate, and
by \eqref{eq:em-missing-info} that is the coordinate whose eigenvalue of $J$ is
pushed closest to one.

\begin{quote}
\emph{Prediction.} $\lambda_{\kappa_4} > \lambda_{\alpha}$, with the gap
governed by fourth- against second-order cumulant damping --- and consequently,
by \eqref{eq:em-settle}, a settling time for the shape that is a large multiple
of the settling time for the correlation.
\end{quote}

\section{Measuring the two rates}\label{sec:em-two-rates}

A dedicated sweep fits the mixture kernel through the channel at
$\nseq = \convsize$, $C = \convcomps$ and the \convdesign{} noise design, over
$\convseeds$ seeds run to $\convcap$ updates, recording the full parameter trace.
Figure~\ref{fig:em-two-rates}(a) plots each coordinate's relative distance to its
own end-of-run value. Both are straight on a log axis over a wide window, which
is \eqref{eq:em-fixedpoint} being obeyed, and the slopes differ sharply.

\begin{center}
\begin{minipage}{\linewidth}\centering
\captionof{table}[The measured EM coordinate contraction factors]{Empirical
\emph{coordinate-wise} contraction
factors $\lambda$ read from the traces of Figure~\ref{fig:em-two-rates}(a)
(median over $\convseeds$ seeds; least-squares slope of
$\log|\theta^{(k)}-\theta^{(\convcap)}|$ over the fitting window, relative
to the end-of-run value, which is itself an estimate of the fixed point). A
coordinate mixes eigenmodes of the ECM map and the excess kurtosis is a
nonlinear functional of the mixture parameters, so these slopes are not
identified with individual Jacobian eigenvalues; no eigenvector analysis is
performed here. The
asymptotic column applies \eqref{eq:em-settle} with the relative tolerance
actually used; the last column is the measured median. The prediction is right
about the ordering and the order of magnitude and over-states the multiple,
for the reason given in the text.}
\label{tab:em-rates}
\small
\setlength{\tabcolsep}{4pt}
\begin{tabular}{@{}lccccc@{}}
\toprule
coordinate & fitting window & $\lambda$ & rel.\ tolerance $\tau$
  & asymptotic $k(\tau,\lambda)$ & measured median \\
\midrule
$\alpha$          & $20$--$150$ & $0.9686$ & $1.2\times10^{-3}$ & $211$  & $\convrhomed$ \\
$\ivar$           & $20$--$150$ & $0.9671$ & ---                & ---    & $\convivarmed$ \\
$\kappa_4$        & $60$--$400$ & $0.9960$ & $6.7\times10^{-3}$ & $1256$ & $\convshapemed$ \\
\midrule
ratio $\kappa_4 / \alpha$ & & $\log$-ratio $8.0$ & & $5.9$ & $\convratio$ \\
\bottomrule
\end{tabular}
\end{minipage}
\end{center}

The prediction of Section~\ref{sec:em-which-slower} is confirmed in direction
and in magnitude of effect. The shape coordinate contracts by a factor
$0.9960$ per update against $0.9686$ for the correlation: in the exponent that
is a factor of $8.0$. Read through \eqref{eq:em-missing-info}, that ordering
is \emph{consistent with} a slow, shape-dominated eigenmode whose missing
information sits far closer to one than the correlation's, exactly the
structure \eqref{eq:em-cumulant-damping} predicts --- consistent with, not a
measurement of: turning the slopes into eigenvalues and missing-information
fractions would require the Jacobian of the implemented map and its
eigenvectors, which are not computed here.
The innovation \emph{variance} settles with the correlation, at
$\lambda = 0.9671$, which is the same statement seen from the other side: it is
a second-order quantity and the channel preserves it.

Where the linear theory is not quantitative is the constant.
Equation~\eqref{eq:em-settle} predicts a settling ratio of $5.9$ and the paired
measurement gives $\convratio$; the asymptotic counts over-predict both
individual settling times by a factor of two to three. The reason is visible in
Figure~\ref{fig:em-two-rates}(a): both curves fall considerably faster in their
first hundred updates than their fitted slopes, so both tolerances are crossed
during the transient rather than in the asymptotic regime that
\eqref{eq:em-fixedpoint} describes. The eigenvalues bound the tail behaviour,
not the crossing of a threshold, and this is the sense in which
Table~\ref{tab:em-rates}'s middle columns are an upper estimate.

\begin{figure}[t]\centering
\includegraphics[width=\linewidth]{fig_em_two_rates.pdf}
\caption[The two convergence rates, and the channel's separate accuracy cost]{
The two effects of Section~\ref{sec:em-two-effects}, measured.
\emph{(a)} Relative distance of each coordinate to its own end-of-run value,
median over $\convseeds$ seeds with the $10$--$90\%$ band; dotted lines are the
fitted contraction rates of Table~\ref{tab:em-rates}. The straightness on a log
axis is \eqref{eq:em-fixedpoint}; the difference in slope is
\eqref{eq:em-missing-info} responding to the cumulant damping
\eqref{eq:em-cumulant-damping}.
\emph{(b)} Settling times per seed, with medians; the correlation and the
innovation variance settle together near $80$ updates and the shape at
$\convshapemed$, with a tail to $\convshapemax$. The dashed line is the fixed
budget of $\convbudget$ updates several experiments used as a default, which
lies below \emph{every} observed settling time of \emph{every} coordinate.
\emph{(c)} The channel's accuracy cost, which is not a rate: relative score
error at $t=0.2$ of a kernel fitted on clean transition pairs against one
fitted on the same chains through the channel, both run to convergence. The
clean arm improves fourfold over the range while the noised arm improves
$2.6$-fold, so the gap \emph{widens} with data rather than closing.}
\label{fig:em-two-rates}
\end{figure}

\section{Settling times and the stopping rule}\label{sec:em-stopping-times}

Figure~\ref{fig:em-two-rates}(b) shows the settling times the rates of
Table~\ref{tab:em-rates} produce. Across $\convseeds$ seeds the correlation
settles at a median of $\convrhomed$ updates and the shape at
$\convshapemed$, the largest observed value being $\convshapemax$. Forming the
ratio \emph{within} each fit and then taking the median gives $\convratio$; the
quotient of the two medians is $\convratioofmedians$. The paired figure is the
one reported, because both coordinates are measured on the same fit.

Two qualifications on what that statistic means. ``Settled'' is trace
stability --- the first update after which a coordinate stays within
$10^{-3}$ of its end-of-run value for $\alpha$, or $2\times10^{-2}$ for the
excess kurtosis --- and not distance to an optimiser; a coordinate can sit
still while the density it belongs to moves along a shallow direction. And the
run has to be long enough for the statistic to mean anything: the sweep was run
at two lengths, and these medians come from the $\convcap$-update shard, the
only one in which no configuration settles near its own cap.

Against them, the fixed budget of $\convbudget$ iterations that several
experiments in Chapter~\ref{ch:em-results} used as a default falls short of the
shape coordinate's median by about a factor of six and leaves
$\convunsettled$ of the sweep's $\convunsettledof$ configurations
--- $\convunsettledpct\%$ --- unsettled in it. That count pools both shards,
which is conservative: censoring can only shorten a measured settling time, so
it cannot inflate the count.

Table~\ref{tab:em-convergence} makes the same point on a single configuration
and adds the endpoints.

\input{sections/tab-convergence}

\noindent The trace shows the coordinates separating exactly as
Table~\ref{tab:em-rates} says they must: the correlation is essentially in
place by thirty updates while the innovation shape is still climbing at a
hundred and twenty. The mixture label is a latent variable of its own, and the
channel compounds the missing information --- the mechanism of
\citet{dempster1977maximum} appearing in the convergence rate, where
\eqref{eq:em-missing-info} places it.

\section{What the channel costs in accuracy at equal data}
\label{sec:em-information}

The endpoints in Table~\ref{tab:em-convergence} are close: the two arms finish
within $0.002$ of each other on the correlation and within a fifth of a unit on
excess kurtosis, both near the generating values. Whether the channel is even
\emph{allowed} to leave the population optimum in place depends on
specification. Under a correctly specified, identifiable model the clean and
the noised likelihoods share their population maximiser, and shared endpoints
are what shrinking $I_{\mathrm{obs}}$ predicts. The fitted family here is a
finite Gaussian mixture that does not contain the Laplace innovation exactly,
and under misspecification the data-processing inequality does not preserve
the \emph{ranking} of candidate parameters, so the clean and noised
pseudo-true optima need not coincide. The close endpoints above are therefore
a \emph{measured approximate invariance} for this family in this regime, not
a theorem --- and either way they are emphatically \emph{not} evidence that
the channel is free. The second effect
lives in the variance, and it has to be measured on the variance.

Figure~\ref{fig:em-two-rates}(c) does so, on the quantity the thesis actually
uses: the relative score error of the fitted kernel at $t = 0.2$, for a kernel
fitted on clean transition pairs against one fitted on the same chains through
the channel, both run to convergence.

\begin{center}
\begin{minipage}{\linewidth}\centering
\captionof{table}[The channel's accuracy cost at equal data]{Relative score
error at $t = 0.2$ of the fitted kernel, clean transition pairs against the
same chains through the channel, $C = 8$, mean over three replicates --- an
exploratory design: three replicates carry no usable uncertainty estimate, so
the row-to-row shape of either curve should not be over-read. The
penalty is real at every size; what the last two columns do and do not
identify is discussed below.}
\label{tab:em-channel-cost}
\small
\begin{tabular}{@{}lccccc@{}}
\toprule
$\nseq$ & $64$ & $128$ & $256$ & $512$ & $1024$ \\
\midrule
clean transition pairs & $0.0229$ & $0.0159$ & $0.0114$ & $0.0071$ & $0.0057$ \\
through the channel    & $0.0451$ & $0.0312$ & $0.0254$ & $0.0179$ & $0.0172$ \\
\midrule
ratio                  & $2.0\times$ & $2.0\times$ & $2.2\times$ & $2.5\times$ & $3.0\times$ \\
\bottomrule
\end{tabular}
\end{minipage}
\end{center}

The penalty is real, and over the range measured it grows: the
clean arm improves fourfold, close to the $\nseq^{-1/2}$ behaviour of
Figure~\ref{fig:em-diagnostics}(c), while the noised arm improves only
$2.6$-fold and is nearly flat between $\nseq = 512$ and $1024$, so the ratio
moves from $2.0\times$ to $3.0\times$.

What that flattening does and does not identify has to be stated carefully,
because the asymptotic theory is unambiguous: under a regular, correctly
specified, identifiable model \emph{both} arms obey
$\sqrt{\nseq}\,(\widehat\theta - \theta^\star) \Rightarrow
\Nd(0, I_{\mathrm{obs}}^{-1})$, so a smaller observed information enlarges
the \emph{constant} of an $\nseq^{-1/2}$ law and cannot, by itself, change
the exponent or produce a plateau. A genuine flattening therefore requires
something else: the finite-mixture approximation floor (the family cannot
represent the Laplace cusp, so the score error cannot fall to zero),
quadrature and stopping tolerances, a possible shift of the noised
pseudo-true target under misspecification, or simply a pre-asymptotic range
--- and three replicates over five sizes cannot separate these. What the
table establishes is a larger finite-data cost through the channel, widening
over the tested range; it does not identify a changed asymptotic rate, and
the flattening is not attributed to Fisher information alone. Reduced
$I_{\mathrm{obs}}$ remains the mechanism for the \emph{level} of the
penalty; the shape of its tail is unresolved here. Equal endpoints
and unequal precision are consistent, and Table~\ref{tab:em-convergence} and
Table~\ref{tab:em-channel-cost} are the two halves of that statement.
